% TOP_PROBLEM: {"id": 30006861, "title": "Irrationality of zeta(5)", "category_id": 1, "category": {"id": 1, "name": "number_theory", "display_name": "Number Theory", "description": "Properties of integers, prime numbers, Diophantine equations.", "slug": "number-theory", "order_index": 1, "created_at": "2026-07-31T15:26:25.670Z"}, "status": "open"}
% FIRST_POSTED: 2026-09-25
% ENTRY_KIND: statement_only
% !TeX program = lualatex
% Definition and sources only; this file is not an LLM solution attempt.
\documentclass[11pt]{article}
\usepackage[margin=25mm]{geometry}
\usepackage{fontspec}
\setmainfont{Latin Modern Roman}
\usepackage{amsmath,amssymb,mathtools}
\usepackage{unicode-math}
\setmathfont{Latin Modern Math}
\usepackage{xurl,hyperref}
\hypersetup{colorlinks=true,urlcolor=blue}
\setcounter{secnumdepth}{0}
\setlength{\parindent}{0pt}
\setlength{\parskip}{5pt}
\setlength{\emergencystretch}{3em}
\begin{document}
\section*{\texorpdfstring{Irrationality of zeta(5)}{Irrationality of zeta(5)}}\label{30006861}
\subsection{Definitions and mathematical statement}
The absolutely convergent series
\[
 \zeta(5)=\sum_{n=1}^{\infty}\frac1{n^5}
\]
defines a positive real number. The problem is to prove $\zeta(5)\notin{\mathbb{Q}}$, equivalently
\[
 \forall a\in{\mathbb{Z}}\ \forall b\in{\mathbb{Z}}_{\ge1},\qquad b\zeta(5)\ne a.
\]
The target singles out this one value, rather than proving that some member of a collection of odd zeta values is irrational. An approximation with many certified decimal digits does not exclude rationality, since the denominator of a hypothetical rational value is unrestricted. Transcendence, meaning that no nonzero polynomial in ${\mathbb{Q}}[X]$ vanishes at this value, is a stronger assertion and is not requested here.
\par\smallskip\textit{Formulation and concept references:} \href{https://www2.mathematik.hu-berlin.de/~kreimer/wp-content/uploads/IrratModuliMotivesv8.pdf}{[S1]}.

\subsection{Short English statement}
Can the sum of the reciprocals of all positive fifth powers be proved not to be a ratio of integers?
\subsection{Sources}
\begin{itemize}
\item[S1]Francis Brown, "Irrationality proofs for zeta values, moduli spaces and dinner parties," author-hosted PDF (v8), www2.mathematik.hu-berlin.de.
\url{https://www2.mathematik.hu-berlin.de/~kreimer/wp-content/uploads/IrratModuliMotivesv8.pdf}.
\textit{Formulation source.}
\end{itemize}
\end{document}
